%%%%%%%%%%%%%%%%%%%%%%%%%%%%%%%%%%%%%%%%%%%%%%%%%%%%%%%%%%%%%%%%%%%%%%%%
% Plantilla TFG/TFM
% Universidad de Murcia. Facultad de Informática
%%%%%%%%%%%%%%%%%%%%%%%%%%%%%%%%%%%%%%%%%%%%%%%%%%%%%%%%%%%%%%%%%%%%%%%%

\chapter{Desarrollo e Implementación}
\label{cap:implementacion}

En este capítulo se detallan los aspectos técnicos de la implementación del sistema. Se exponen las decisiones de programación, las herramientas utilizadas y fragmentos de código clave que ilustran la lógica interna del chatbot.

\section{Entorno de Desarrollo y Tecnologías}

El sistema ha sido desarrollado íntegramente en Python, aprovechando su extenso ecosistema para el procesamiento de datos y la Inteligencia Artificial. Entre las librerías principales destacan:
\begin{itemize}
    \item \textbf{LangChain:} Para la orquestación del LLM y la gestión del bucle del agente.
    \item \textbf{RDFLib:} Para la carga, manipulación y consulta en memoria del Grafo de Conocimiento (\texttt{.ttl}).
    \item \textbf{Streamlit:} Para la rápida construcción de la interfaz gráfica de usuario.
    \item \textbf{SPARQLWrapper:} Para realizar consultas remotas a Wikidata durante el proceso de extracción de datos y desambiguación de entidades.
\end{itemize}

El despliegue del sistema se ha contenerizado utilizando \textbf{Docker}, garantizando así que el entorno de ejecución sea reproducible en cualquier máquina sin problemas de dependencias.

\section{Extracción de Datos y Grafo Local}

Para evitar depender de llamadas externas costosas durante la interacción con los usuarios, se diseñó el proceso (\texttt{generar\_datos\_cine.py}) que extrae la información desde Wikidata y la guarda en el archivo \texttt{datos\_cine\_seguro.ttl}, que funciona como base de datos local.

Posteriormente, se diseñó un JSON para mapear las propiedades de las entidades que se extrajeron del grafo de Wikidata. Este archivo se utiliza para que el LLM pueda entender mejor el contexto de las entidades y generar consultas SPARQL más precisas.

La estructura del código permite añadir entidades relevantes al grafo local. Por ejemplo, al extraer una película, también se extraen los nodos de su director y actores, conformando un subgrafo fuertemente conectado, ideal para responder a las consultas SPARQL posteriores.

\section{El Agente Conversacional y el Control del LLM}

El componente principal, ubicado en \texttt{chatLocal.py}, inicializa el motor LangChain. El aspecto más crítico de la implementación es el diseño de los \textit{Prompts} del sistema. 

Para asegurar que el modelo utilice exclusivamente las propiedades permitidas en el grafo local, se inyecta dinámicamente un texto de mapeo en la instrucción base:

\begin{lstlisting}[style=Python-color, caption={Inyección de mapeo de propiedades en el prompt}, label=lst:prompt]
with open("mapeo_propiedades.json", "r", encoding="utf-8") as f:
    mapeoProp = json.load(f)
    textoMapeo = json.dumps(mapeoProp, indent=2)

system_template = f"""
Eres un agente experto en SPARQL y bases de datos RDF del dominio de cine.
Esquema de propiedades disponibles:
{textoMapeo}
Objetivo: 
Genera una consulta SPARQL válida para responder a la consulta '{consulta}'
REGLAS:
    1. La entidad principal es: wd:{id_entidad}
    2. Usa SOLO las propiedades del esquema de arriba.
    3. Las propiedades DEBEN llevar el prefijo 'wdt:'. Si usas la propiedad P178, escribe wdt:P178.
    ...
"""
\end{lstlisting}

\section{Bucle de Reflexión y Autocorrección}

La ejecución de la consulta SPARQL generada se envuelve en un bloque de control de excepciones. Si \textit{RDFLib} arroja un error de sintaxis al evaluar la consulta, el sistema captura dicho error y realiza una llamada iterativa al LLM para que se corrija a sí mismo:

\begin{lstlisting}[style=Python-color, caption={Bucle de reflexión ante errores de sintaxis}, label=lst:reflexion]
try:
    resultados = g.query(query_sparql)
except Exception as e:
    # Si falla, informamos al agente de su propio error
    mensaje_error = f"Error sintáctico: {str(e)}. Reescribe el SPARQL."
    nuevo_sparql = llm.invoke(prompt_reflexion + mensaje_error)
    resultados = g.query(nuevo_sparql)
\end{lstlisting}

Este enfoque incrementó drásticamente la robustez del sistema, permitiéndole recuperarse de fallos comunes como puntos finales olvidados o prefijos \texttt{wdt} mal estructurados, sin mostrar errores en bruto al usuario final.

\section{Interfaz Gráfica}

La interfaz, construida con \textit{Streamlit} (\texttt{interfaz.py}), mantiene el estado de la conversación y proporciona una experiencia moderna mediante técnicas de diseño \textit{Glassmorphism}. Destaca la inclusión de un \textbf{panel desplegable de estado} donde el usuario puede visualizar en tiempo real los pensamientos internos del agente: qué entidad extrae, qué consulta SPARQL ha generado y si ha necesitado autocorregirse. Esto dota al sistema de una robusta capa de explicabilidad (\textit{Explainable AI}), diferenciándolo de los chatbots tipo "caja negra" convencionales.
